\chapter{Progettazione dell'Applicazione Software}

\section{Architettura Software Complessiva}
L'applicazione software, denominata \textbf{Wedding Planner Studio Desktop}, è stata realizzata in linguaggio \textbf{Java 21 LTS} utilizzando lo strumento di automazione \textbf{Gradle 9.2} e il driver ufficiale \textbf{MySQL Connector/J 8.3.0}.

L'architettura del sistema adotta una strutturazione a livelli (\textit{3-Tier Layered Architecture}) basata sul pattern \textbf{Data Access Object (DAO)}:

\begin{enumerate}
    \item \textbf{Presentation Layer (\texttt{db\_lab.ui}):} Interfaccia grafica desktop sviluppata con la libreria \textbf{Java Swing}, integrata con il Look \& Feel nativo del sistema operativo per un'esperienza d'uso fluida ed ergonomica. La finestra principale (\texttt{MainWindow.java}) è organizzata in un \texttt{JTabbedPane} modulare a 6 schede.
    \item \textbf{Data Access Layer (\texttt{db\_lab.dao}):} Modulo responsabile dell'incapsulamento dell'accesso alla base di dati, della composizione dei \texttt{PreparedStatement} per la prevenzione totale da attacchi \textit{SQL Injection} e della conversione dei \texttt{ResultSet} in strutture dati tipizzate.
    \item \textbf{Persistence Layer (MySQL 8.0):} Storage engine InnoDB locale su porta 3306, con gestione transazionale e trigger attivi.
\end{enumerate}

\section{Organizzazione dei Moduli e Package}
I sorgenti applicativi sono organizzati nei seguenti package:
\begin{itemize}[noitemsep]
    \item \texttt{db\_lab.App}: Entry point dell'applicazione; inizializza il Look \& Feel e stabilisce la connessione JDBC tramite credenziali dinamiche.
    \item \texttt{db\_lab.ui.MainWindow}: Controller e vista unificata della GUI desktop; gestisce eventi, rendering tabellare e finestre modali.
    \item \texttt{db\_lab.data.DAOUtils}: Utility factory per connessioni JDBC e metodo generico \texttt{executeGenericQuery(...)}.
    \item \texttt{db\_lab.data.QueryResult}: Record immutabile Java 21 che trasporta metadati delle colonne, righe di dati, conteggio record e tempo di esecuzione in millisecondi, con serializzazione diretta in vettori compatibili con \texttt{DefaultTableModel}.
    \item \texttt{db\_lab.dao.StatisticheDAO}: Motore dedicato all'esecuzione e metadati delle 8 query analitiche (\texttt{STAT-01} .. \texttt{STAT-08}).
    \item \texttt{db\_lab.dao.EventoDAO}: Consultazione anagrafica matrimoni, budget allocato, planner assegnati e totale contratti.
    \item \texttt{db\_lab.dao.TavoliDAO}: Navigazione tavoli della sala, elenco invitati per matrimonio e metodo transazionale di assegnazione posto \texttt{assegnaTavolo(...)}.
    \item \texttt{db\_lab.dao.FornitoreDAO}: Anagrafica fornitori per categoria merceologica e contratti per matrimonio.
    \item \texttt{db\_lab.dao.ListaNozzeDAO}: Lettura coordinate bancarie, regali, quote ricevute e inserimento di nuove donazioni.
\end{itemize}

\section{Descrizione Dettagliata delle Schede dell'Applicazione}

\subsection{Scheda 1: Statistiche Direzionali (STAT-01 .. STAT-08)}
Costituisce il cruscotto direzionale avanzato per la direzione e i planner. Consente di:
\begin{itemize}[noitemsep]
    \item Selezionare tramite menu a discesa una qualsiasi delle 8 interrogazioni statistiche previste dalle specifiche.
    \item Visualizzare il titolo, la descrizione estesa dell'obiettivo aziendale e, opzionalmente, il codice SQL eseguito.
    \item Premere il pulsante \textbf{«⚡ Esegui Statistica»} per popolare istantaneamente la tabella risultati.
    \item Ordinare dinamicamente le righe cliccando sulle intestazioni di colonna (\texttt{TableRowSorter}).
    \item Monitorare nella status bar inferiore il tempo di esecuzione preciso (tipicamente inferiore a 15 ms) e il numero di tuple restituite.
\end{itemize}

\subsection{Scheda 2: Panoramica Eventi \& Matrimoni}
Offre una vista a volo d'uccello di tutti i matrimoni registrati a sistema:
\begin{itemize}[noitemsep]
    \item Dati esposti: Coppia committente, data fissata, stato avanzamento (\texttt{in\_pianificazione}, \texttt{confermato}, \texttt{completato}), budget effettivo, invitati previsti, tipo di rito, tema scenografico, wedding planner responsabile e totale dei contratti fornitore impegnati.
    \item Pulsante di aggiornamento rapido in tempo reale.
\end{itemize}

\subsection{Scheda 3: Table Planner \& Gestione Ospiti (RSVP)}
Progettata specificamente per l'allestimento della sala e la sicurezza alimentare degli ospiti:
\begin{itemize}[noitemsep]
    \item \textbf{Pannello sinistro (Tavoli):} Elenco dei tavoli con nome/numero, capienza massima, posti occupati, posti liberi e stato di saturazione (\texttt{DISPONIBILE}, \texttt{COMPLETO}, \texttt{VUOTO}).
    \item \textbf{Pannello destro (Invitati):} Elenco partecipanti con codice fiscale, nominativo, recapiti, stato di conferma (\texttt{CONFERMATO}, \texttt{IN\_ATTESA}, \texttt{RIFIUTATO}), tipologia di menù e indicazione chiara di allergie ed intolleranze.
    \item \textbf{Pannello inferiore di Assegnazione Posto:} Selezionando un invitato e scegliendo un tavolo, il pulsante \textbf{«Assegna / Sposta Tavolo»} esegue la transazione validando a livello software e DB i vincoli di business (impedendo ad esempio l'assegnazione a chi ha stato \texttt{IN\_ATTESA} o il superamento della capienza massima).
\end{itemize}

\subsection{Scheda 4: Fornitori \& Contratti}
Consente la gestione del catalogo partner:
\begin{itemize}[noitemsep]
    \item Filtro dinamico per settore merceologico (\textit{Catering, Fotografia, Fiori, Musica, Abiti, Tutte}).
    \item Monitoraggio del numero di contratti stipulati e del fatturato complessivo intermediato.
    \item Nel pannello inferiore, selezione del matrimonio per consultare il dettaglio economico dei singoli contratti: fornitore, data firma, importo concordato, acconto già corrisposto e saldo ancora da saldare.
\end{itemize}

\subsection{Scheda 5: Lista Nozze \& Donazioni}
Gestione trasparente dei contributi economici degli invitati:
\begin{itemize}[noitemsep]
    \item Banner superiore con coordinate bancarie (IBAN, Banca, Intestatario) e progresso percentuale del totale raccolto rispetto al totale desiderato dalla coppia.
    \item Tabella articoli desiderati con target di prezzo e percentuale di copertura raggiunta.
    \item Tabella storico quote versate con data, nome donatore, metodo di pagamento e messaggio di auguri.
    \item Pulsante \textbf{«💳 Nuova Quota Regalo»}: apre una finestra di dialogo modale che consente di registrare una donazione live nel database, validando il donatore ed azionando il trigger che aggiorna istantaneamente il saldo della lista nozze.
\end{itemize}

\subsection{Scheda 6: Console SQL Interattiva}
Inserita per agevolare le fasi di collaudo e la discussione orale d'esame con i docenti:
\begin{itemize}[noitemsep]
    \item Area di testo per inserire qualsiasi query libera \texttt{SELECT}.
    \item Pulsanti di accesso rapido per interrogazioni frequenti (Elenco Coppie, Attività Checklist con Ritardi, Fornitori attivi, Recensioni).
    \item Rendering automatico e dinamico della struttura delle colonne e dei dati restituiti.
\end{itemize}

\section{Verifica e Test Automatizzati del Software}
L'affidabilità del sistema è stata garantita tramite una suite di test automatici scritta con \textbf{JUnit 5} e \textbf{AssertJ} nella classe \texttt{WeddingPlannerTest.java}:
\begin{itemize}[noitemsep]
    \item \textbf{Test di Connessione:} Accerta la raggiungibilità del database \texttt{wedding\_planner} su porta 3306.
    \item \textbf{Test Parametrizzato Query Analitiche:} Esegue a rotazione tutte le 8 query analitiche (\texttt{STAT-01} fino a \texttt{STAT-08}), verificando che il set di risultati sia non nullo, che il numero di tuple restituite sia coerente con il dataset e che i tempi di esecuzione siano nell'ordine dei millisecondi.
    \item \textbf{Test dei DAO di Business:} Verifica la correttezza delle query di \texttt{EventoDAO}, \texttt{TavoliDAO}, \texttt{FornitoreDAO} e \texttt{ListaNozzeDAO}.
\end{itemize}

L'esecuzione della suite di collaudo:
\begin{lstlisting}[language=bash]
./gradlew test
\end{lstlisting}
ha registrato il completamento con successo di tutti i \textbf{17 test} eseguiti con \textbf{zero errori o fallimenti} (\textit{BUILD SUCCESSFUL}).
